\subsection{ASYMPTOTIC EXPANSION OF THE PARTIAL SUMS OF A SERIES}

For \(x\) real tending to \(+\infty\) let \(\mathcal{E}\) be a comparison scale formed by functions each of which is defined on a whole interval \([x_0, +\infty[\) (depending on the function) and is \(\geqslant 0\) on this interval. Let \((\mathbf{u}_n)\) be a series whose terms belong to a complete normed space \(E\) , such that \(\mathbf{u}_n\) admits an asymptotic expansion to precision \(g_\alpha\) with respect to the scale \(\mathcal{E}'\) of restrictions to \(\mathbf{N}\) of the functions in \(\mathcal{E}\) :

\[
\mathbf {u} _ {n} = \sum_ {\lambda \leqslant \alpha} \mathbf {a} _ {\lambda} g _ {\lambda} (n) + \mathbf {r} _ {\alpha} (n).
\]

Suppose that every partial sum \(\sum_{m=1}^{n} g(m)\) , where \(g \in \mathcal{E}\) , admits an asymptotic expansion relative to \(\mathcal{E}'\) . One can then obtain an asymptotic expansion of the \(\mathbf{s}_n = \sum_{m=1}^{n} \mathbf{u}_m\) with respect to \(\mathcal{E}'\) ; again we distinguish two cases:

\(1^{\circ}\sum_{n = 1}^{\infty}g_{\alpha}(n) = +\infty .\) Then (V,p.237,prop.2) one has \(\sum_{m = 1}^{n}\mathbf{r}_{\alpha}(m)\ll \sum_{m = 1}^{n}g_{\alpha}(m);\)

by hypothesis one can obtain an asymptotic expansion of

\[
\sum_ {\lambda \leqslant \alpha} \mathbf {a} _ {\lambda} \left(\sum_ {m = 1} ^ {n} g (m)\right)
\]

(V, p.~222) to a certain precision \(g_{\sigma}\) ; if \(cg_{\sigma}(n)\) is the principal part of \(\sum_{m=1}^{n} g_{\alpha}(m)\) , one will have an asymptotic expansion of \(\mathbf{s}_n\) to precision \(g_{\min(\rho, \sigma)}\) .

\(2^{\circ}\sum_{n = 1}^{\infty}g_{\alpha}(n)\) converges; then let \(\beta\) be the smallest of the indices \(\lambda \leqslant \alpha\) such that \(\mathbf{a}_{\lambda}\neq 0\) and such that \(\sum_{n = 1}^{\infty}g_{\lambda}(n)\) converges; the series

\[
\mathbf {C} = \sum_ {n = 1} ^ {\infty} \left(\mathbf {u} _ {n} - \sum_ {\lambda <   \beta} \mathbf {a} _ {\lambda} g _ {\lambda} (n)\right)
\]

then converges absolutely, and one can write

\[
\mathbf {s} _ {n} = \sum_ {\lambda <   \beta} \mathbf {a} _ {\lambda} \left(\sum_ {m = 1} ^ {n} g _ {\lambda} (m)\right) + \mathbf {C} - \sum_ {\beta \leqslant \lambda \leqslant \alpha} \mathbf {a} _ {\lambda} \left(\sum_ {m = n + 1} ^ {\infty} g _ {\lambda} (m)\right) - \sum_ {m = n + 1} ^ {\infty} \mathbf {r} _ {\alpha} (m).
\]

Further, \(\sum_{m=n+1}^{\infty}\mathbf{r}_{\alpha}(m)\ll\sum_{m=n+1}^{\infty}g_{\alpha}(m)\) ; if \(cg_{\sigma}(n)\) is the principal part of \(\sum_{m=n+1}^{\infty}g_{\alpha}(m)\) , and if one has an asymptotic expansion of

\[
\sum_ {\lambda <   \beta} \mathbf {a} _ {\lambda} \left(\sum_ {m = 1} ^ {n} g _ {\lambda} (m)\right) + \mathbf {C} - \sum_ {\beta \leqslant \lambda \leqslant \alpha} \mathbf {a} _ {\lambda} \left(\sum_ {m = n + 1} ^ {\infty} g _ {\lambda} (m)\right)
\]

to precision \(g_{\rho}\) one thus obtains an asymptotic expansion of \(s_{n}\) to precision \(g_{\min(\rho,\sigma)}\) . One is thus led to the particular case of series \((g(n))\) where \(g \in E\) . We shall see how, subject to certain conditions, one can straight away obtain a principal part of \(s_{n} = \sum_{m=1}^{n} g(m)\) (when \(\sum_{n=1}^{\infty} g(n) = +\infty\) ) or of \(r_{n} = \sum_{m=n+1}^{\infty} g(m)\) (when \(\sum_{n=1}^{n} g(n) < +\infty\) ).

PROPOSITION 6. Let \(g\) be a real function, \(>0\) and monotone, defined on an interval \([x_0, +\infty[\) (where \(x_0 \leqslant 1\) ), and such that \(\log g\) and \(x\) are comparable to order 1.

\(1^{\circ}\) If \(g\) is of infinite order relative to \(e^{x}\) one has

\[
s _ {n} = \sum_ {m = 1} ^ {n} g (m) \sim g (n) \quad i f \quad \sum_ {n = 1} ^ {\infty} g (n) = + \infty ;\tag{1}
\]

\[
r _ {n} = \sum_ {m = n + 1} ^ {\infty} g (m) \sim g (n + 1) \quad i f \quad \sum_ {n = 1} ^ {\infty} g (n) <   + \infty .\tag{2}
\]

\(2^{\circ}\) If \(g\) is of finite order \(\mu\) with respect to \(e^{x}\) one has

\[
s _ {n} = \sum_ {m = 1} ^ {n} g (m) \sim \frac {\mu}{1 - e ^ {- \mu}} \int_ {x _ {0}} ^ {n} g (t) d t \quad i f \quad \sum_ {n = 1} ^ {\infty} g (n) = + \infty ;\tag{3}
\]

\[
r _ {n} = \sum_ {m = n + 1} ^ {\infty} g (m) \sim \frac {\mu}{1 - e ^ {- \mu}} \int_ {n} ^ {\infty} g (t) d t \quad i f \quad \sum_ {n = 1} ^ {\infty} g (n) <   + \infty\tag{4}
\]

(the number \(\frac{\mu}{1 - e^{-\mu}}\) is to be replaced by 1 in (3) and (4) when \(\mu = 0\) ).

\(1^{\circ}\) If \(g\) is of order \(+\infty\) relative to \(e^{x}\) one has \(\log g \gg x\) whence \(g' / g \gg 1\) , or \(g' \gg g\) , by the hypothesis; \(g\) is therefore increasing and tends to \(+\infty\) with \(x\) , whence \(\sum_{n=1}^{\infty} g(n) = +\infty\) . If \(u\) is the step function associated with the series \((g(n))\) (V, p.~237), one has \(u(x) \leqslant g(x)\) starting from a certain value of \(x\) , so \(u \preccurlyeq g\) and consequently

\[
s _ {n - 1} = \int_ {1} ^ {n} u (t) d t \preccurlyeq \int_ {1} ^ {n} g (t) d t \prec \prec \int_ {1} ^ {n} g ^ {\prime} (t) d t \sim g (n);
\]

since \(s_n = s_{n-1} + g(n)\) one has \(s_n \sim g(n)\) . The proof is similar when \(g\) is of order \(-\infty\) relative to \(e^x\) ; we thus obtain formula (2).

\(2^{\circ}\) If \(g\) is of finite order \(\mu\) relative to \(e^{x}\) one can write \(g(x) = e^{\mu x}h(x)\) where \(h\) is of order 0 relative to \(e^{x}\) ; further, by hypothesis, \(\log g \sim \mu x\) for \(\mu \neq 0\) ( \(\log g \prec x\) for \(\mu = 0\) ) implies \(h' \prec h\) . Suppose first that \(\sum_{n=1}^{\infty} g(n) = +\infty\) (which implies that \(\mu \geqslant 0\) ; the converse always holds if \(\mu > 0\) , since then \(g(x)\) tends to \(+\infty\) with \(x\) ); let us evaluate the principal part of \(\int_{n-1}^{n} g(t) dt\) . One can write

\[
\begin{array}{l} \int_ {n - 1} ^ {n} g (t) d t = \int_ {n - 1} ^ {n} e ^ {\mu t} h (t) d t \\ \qquad = h (n) \int_ {n - 1} ^ {n} e ^ {\mu t} d t + \int_ {n - 1} ^ {n} e ^ {\mu t} \big (h (t) - h (n) \big) d t \\ \qquad = \frac {1 - e ^ {- \mu}}{\mu} g (n) + \int_ {n - 1} ^ {n} e ^ {\mu t} \big (h (t) - h (n) \big) d t. \end{array}
\]

Now, the relation \(h^{\prime} \ll h\) implies that for every \(\varepsilon > 0\) there is an \(n_{0}\) such that the relation \(x \geqslant n_{0}\) implies that \(\left|h'(x)/h(x)\right| \leqslant \varepsilon\) ; one deduces from the mean value theorem that \(-\varepsilon \leqslant \log |h(t)/h(n)| \leqslant \varepsilon\) , for \(n - 1 \leqslant t \leqslant n\) , if \(n \geqslant n_{0}\) , whence

\[
| h (t) - h (n) | \leqslant \left(e ^ {\varepsilon} - 1\right) h (n)
\]

and consequently

\[
\left| \int_ {n - 1} ^ {n} e ^ {\mu t} (h (t) - h (n)) d t \right| \leqslant (e ^ {\varepsilon} - 1) e ^ {\mu n} h (n) = (e ^ {\varepsilon} - 1) g (n)
\]

since \(e^{\mu t}\) is increasing. Since \(e^{\varepsilon} - 1\) becomes arbitrarily small with \(\varepsilon\) one sees that one can write

\[
\int_ {n - 1} ^ {n} g (t) d t = \frac {1 - e ^ {- \mu}}{\mu} g (n) + o (g (n))
\]

\(\left(\frac{1 - e^{-\mu}}{\mu}\right.\) being replaced by 1 when \(\mu = 0\) . The proposition is then a consequence of prop. 2 of V, p.~237. One argues similarly when \(\sum_{n=1}^{\infty} g(n)\) is finite.

By applying prop. 6 of V, p.~239, repeatedly one can then sometimes obtain an asymptotic expansion for \(s_{n} = \sum_{m=1}^{n} g(m)\) . First suppose that g is of order \(+\infty\) relative to \(e^{x}\) ; for every fixed value of p one can write, by prop. 6,

\[
s _ {n} = g (n) + g (n - 1) + \dots + g (n - p) + o (g (n - p))
\]

and it suffices to expand (relative to \(E'\) ) each of the functions \(g(n-k)\) ( \(0 \leqslant k \leqslant p\) ), limiting the precision of the expansions to the principal part of \(g(n-p)\) , to obtain an expansion for the \(s_{n}\) .

Example. Let \(g(x) = x^{x} = \exp (x\log x)\) , of order \(+\infty\) relative to \(e^x\) . Taking \(p = 2\) one has

\[
(n - 1) \log (n - 1) = (n - 1) \log n - 1 + \frac {1}{2 n} + o \left(\frac {1}{n}\right)
\]

whence (V, p.~225)

\[
(n - 1) ^ {n - 1} = \frac {1}{e} n ^ {n - 1} + \frac {1}{2 e} n ^ {n - 2} + o _ {1} (n ^ {n - 2})
\]

and similarly

\[
(n - 2) ^ {n - 2} = \frac {1}{e ^ {2}} n ^ {n - 2} + o _ {2} (n ^ {n - 2});
\]

consequently

\[
s _ {n} = n ^ {n} + \frac {1}{e} n ^ {n - 1} + \left(\frac {1}{2 e} + \frac {1}{e ^ {2}}\right) n ^ {n - 2} + o _ {3} (n ^ {n - 2}).
\]

One proceeds similarly (for \(r_n\) ) when \(g\) is of order \(-\infty\) relative to \(e^x\) .

Now if \(g\) is of finite order \(\mu\) relative to \(e^x\) , and if, for example, \(\sum_{n=1}^{\infty} g(n) = +\infty\) , one can write

\[
s _ {n} = \frac {\mu}{1 - e ^ {- \mu}} \int_ {1} ^ {n} g (t) d t + \sum_ {m = 1} ^ {n} f _ {1} (m)
\]

where \(f_{1}(n) = g(n) - \frac{\mu}{1 - e^{-\mu}}\int_{1}^{n}g(t)dt \ll g(n)\) by prop. 6 of V, p.~239. If one has a principal part \(cg_{1}(n)\) of \(f_{1}(n)\) relative to \(\mathcal{E}'\) , and if one can again apply prop. 6 to the function \(g_{1}\) one will obtain a primitive equivalent to \(\sum_{m=1}^{n}f_{1}(m)\) if \(\sum_{n=1}^{\infty}g_{1}(n) = +\infty\) , and equivalent to \(\sum_{m=n+1}^{\infty}f_{1}(m)\) in the opposite case (in the latter case one writes \(\sum_{m=1}^{n}f_{1}(m) = C - \sum_{m=n+1}^{\infty}f_{1}(m)\) , with \(C = \sum_{n=1}^{\infty}f_{1}(n)\) ).

Step by step one can thus eventually obtain an expression for \(s_{n}\) as the sum of a certain number of primitives each of which is negligible with respect to the previous, of a term remaining negligible relative to the last primitive written, and finally a constant (the case where the remainder term tends to 0). It then remains to expand each of the primitives obtained with respect to \(E'\) (cf.~V, p.~235).

Example. Let \(g(n) = \frac{1}{n}\) ; then

\[
s _ {n} = \sum_ {m = 1} ^ {n} \frac {1}{m} \sim \int_ {1} ^ {n} \frac {d t}{t} = \log n
\]

then

\[
\frac {1}{n} - (\log n - \log (n - 1)) \sim - \frac {1}{2 n ^ {2}}
\]

whence

\[
s _ {n} = \log n + \gamma + \frac {1}{2 n} + o \left(\frac {1}{n}\right).
\]

The constant \(\gamma\) which appears in this formula plays an important rôle in Analysis (cf.~chap.~VI and VII); it is known as Euler's constant; one has

\[
\gamma = 0. 5 7 7   2 1 5   6 6 4 \dots
\]

to within \(1/10^{9}\) .

We shall see in VI, p.~288, how the Euler-Maclaurin summation formula gives an asymptotic expansion of arbitrary order for \(s_n\) (or for \(r_n\) ) in the most important cases.

